% Part: first-order-logic
% Chapter: beyond
% Section: intuitionistic-logic

\documentclass[../../../include/open-logic-section]{subfiles}

\begin{document}

\olfileid{fol}{byd}{il}

\olsection{شهودي منطق}

د دويمې او لوړې درجې د منطق پر خلاف، د لومړۍ درجې شهودي منطق د
کلاسیکې بڼې يو محدود شوے نظام دے چې د لا ``تعميري'' ډول استدلال د مدل
کولو دپاره جوړ شوے دے۔ لاندې بېلګې ښايي ځينې بنسټيزې انګېزې روښانه کړي۔

فرض کړئ يوه ورځ يو څوک درته راشي او اعلان وکړي چې طبيعي عدد~$x$ ئې
ټاکلے دے او دا خاصيت لري: که $x$ اولي وي، د ريمان فرضيه رښتونې ده؛
او که $x$ مرکب وي، د ريمان فرضيه دروغ ده۔ دا خو لوے زېرے دے!{} د
ريمان د فرضيې رښتياوالے يا دروغوالے د رياضياتو له لويو ناحل شويو
پوښتنو څخه دے، او داسې ښکاري چې دې کس مسئله محاسبې ته راکمه کړې،
يعنې يوازې بايد معلومه شي چې يو ځانګړے عدد اولي دے که نۀ۔

د~$x$ جادويي قيمت څو دے؟ هغه ئې داسې بيانوي: $x$ هغه طبيعي عدد دے
چې که د ريمان فرضيه رښتونې وي له~$7$ سره برابر دے، او که نۀ،
له~$9$ سره۔

تاسو په غوسه خپلې پيسې بېرته غواړئ۔ له کلاسیک نظره پاسنے بيان په
رښتيا د~$x$ يو يوازينے قيمت ټاکي؛ خو څه چې تاسو ئې په اصل کښې غواړئ
د~$x$ داسې قيمت دے چې \emph{په څرګند ډول} درکړل شوے وي۔

د يوې بلې، ښايي لږې مصنوعي، بېلګې دپاره لاندې پوښتنه وګورئ۔ پوهېږو
چې يو غير ناطق عدد ناطق توان ته پورته کېدے او ناطقه پايله ورکولے شي؛
د بېلګې په توګه، $\sqrt{2}^2 = 2$۔ دومره څرګنده نۀ ده چې يو غير ناطق
عدد \emph{غير ناطق} توان ته پورته شي او بيا هم ناطقه پايله ورکړي که
نۀ۔ لاندې قضيه مثبت ځواب ورکوي:

\begin{thm}
داسې غير ناطق عددونه $a$ او $b$ شته چې $a^b$ ناطق وي۔
\end{thm}

\begin{proof}
$\sqrt{2}^{\sqrt{2}}$ وګورئ۔ که دا ناطق وي، کار بشپړ دے: کولے شو
$a = b = \sqrt{2}$ ونيسو۔ که نۀ وي، دا غير ناطق دے۔ بيا لرو:
\[
(\sqrt{2}^{\sqrt{2}})^{\sqrt{2}} = \sqrt{2}^{\sqrt{2} \cdot
  \sqrt{2}} = \sqrt{2}^2 = 2,
\]
چې خامخا ناطق دے۔ نو په دې حالت کښې $a$ له
$\sqrt{2}^{\sqrt{2}}$ سره او $b$ له~$\sqrt 2$ سره برابر ونيسئ۔
\end{proof}

آيا دا يو معتبر ثبوت دے؟ زياتره رياضيپوهان ئې معتبر ګڼي۔ خو دلته بيا
هم څه ناڅه نۀ قانع کوونکې خبره شته: د حقيقي عددونو د يوې جوړې شتون
مو په يوه خاصيت سره ثابت کړے، بې له دې چې ويلے شو دا \emph{کومه}
جوړه ده۔ همدا پايله داسې هم ثابتېدے شي چې جوړه~$a$، $b$ پخپله په
ثبوت کښې \emph{ورکړل شوې} وي: $a = \sqrt{3}$ او $b = \log_3 4$
ونيسئ۔ بيا
\[
a^b = \sqrt{3}^{\log_3 4} = 3^{1/2 \cdot \log_3 4} = (3^{\log_3
  4})^{1/2} = 4^{1/2}= 2,
\]
ځکه چې $3^{\log_3 x} = x$۔

شهودي منطق د داسې استدلال د مدل کولو دپاره جوړ شوے چې د لومړي ثبوت
په څېر ګامونه نۀ مني۔ د~$!A(x)$ پوره کوونکي~$x$ د شتون ثابتول په دې
معنا دي چې يو ځانګړے~$x$ او د هغه له خوا د~$!A$ د پوره کولو ثبوت
ورکړئ، لکه په دويم ثبوت کښې۔ دا ثابتول چې $!A$ يا $!B$ صدق کوي،
غواړي چې له دواړو څخه يو په رښتيا ثابت کړے شئ۔

په صوري ډول، د لومړۍ درجې شهودي منطق هغه وخت ترلاسه کېږي چې د لومړۍ
درجې د منطق !!a{derivation} نظام په يوه ټاکلې طريقه محدود کړئ۔ همداسې
د دويمې يا لوړې درجې منطق شهودي بڼې هم شته۔ له رياضيکي نظره دا يوازې
صوري استنتاجي نظامونه دي، خو لکه مخکې چې وويل شول، موخه ئې د رياضيکي
استدلال يو ډول مدل کول دي۔ څوک دا هغه استدلال ګڼلے شي چې د رياضياتو
يو ځانګړے فلسفي دريځ ئې روا بولي (لکه د بروور شهوديت)؛ څوک ئې د
رياضيکي استدلال يو لا ``ملموس'' او قانع کوونکے ډول ګڼلے شي (د بيشپ د
تعميريت په دود)؛ او دا بحث هم کېدے شي چې صوري بيان غيرصوري انګېزه په
رښتيا رانغاړي که نۀ۔ خو هر فلسفي دريځ چې ولرو، شهودي منطق د يوه صوري
وړاندې شوي منطق په توګه مطالعه کولے شو؛ او د هر دليل له مخې، ډېر
رياضيکي منطقپوهان دا مطالعه په زړه پورې ګڼي۔

د شهودي تړونکو يو غيرصوري تعميري تفسير شته چې په عام ډول د بي‌اېچ‌کې
تفسير بلل کېږي (د بروور، هايتنګ او کولموګوروف په نومونو نومول شوے)۔
دا داسې دے: د~$!A \land !B$ ثبوت د~$!A$ له يوه ثبوت او د~$!B$ له
يوه ثبوت څخه جوړه وي؛ د~$!A \lor !B$ ثبوت يا د~$!A$ ثبوت وي يا
د~$!B$، او څرګند معلومات لرو چې کوم يو دے؛ د~$!A \lif !B$ ثبوت يوه
کړنلاره وي چې د~$!A$ ثبوت د~$!B$ په ثبوت بدلوي؛ د
$\lforall[x][!A(x)]$ ثبوت يوه کړنلاره وي چې د~$x$ د هر قيمت دپاره د
$!A(x)$ ثبوت بېرته ورکوي؛ او د~$\lexists[x][!A(x)]$ ثبوت د~$x$ له
يوه قيمت او دې ثبوت څخه جوړ وي چې هغه قيمت~$!A$ پوره کوي۔ دا تفسير
په هغو محاسبوي اصطلاحاتو هم بيانېدے شي چې د ``کري--هاورډ
آيزومورفيزم'' يا ``د ټايپونو په توګه د !!{formula}s بېلګه'' بلل
کېږي: !!a{formula} د معلوماتو يو ټاکلے ټايپ مشخصوونکے وګڼئ، او ثبوتونه
د همدې معلوماتي ټايپونو محاسبوي څيزونه وګڼئ چې راښيي اړوند
!!{formula} رښتونے دے۔

شهودي منطق ډېر ځله کلاسیک منطق ``د خارج وسط له اصل پرته'' ګڼل کېږي۔
لاندې قضيه دا خبره لا دقيقه کوي۔
\begin{thm}
په شهودي ډول لاندې بديهي شېمې يو له بل سره برابرې دي:
\begin{enumerate}
\item $(\lnot !A \lif \lfalse) \lif !A$.
\item $!A \lor \lnot !A$
\item $\lnot \lnot !A \lif !A$
\end{enumerate}
\end{thm}
له دوو نورو څخه د هرې شېمې بېلګې ترلاسه کول په شهودي منطق کښې يو ښه
تمرين دے۔

د شهودي بياني منطق لومړني استنتاجي نظامونه، چې د بروور د شهوديت د
صوري بڼو په توګه وړاندې شوي وو، کولموګوروف، ګليوېنکو او هايتنګ په
خپلواکه توګه جوړ کړل۔ د لومړۍ درجې شهودي منطق (او د شهودي رياضياتو
د ځينو برخو) لومړۍ صوري بڼه هايتنګ وړاندې کړه۔ که څه هم يو شمېر په
کلاسیک ډول معتبرې شېمې په شهودي ډول معتبرې نۀ دي، ډېرې نورې معتبرې
دي۔

\emph{دوه‌ګونې نفي ژباړه} د کلاسیک او شهودي منطق تر منځ يوه مهمه
اړيکه بيانوي۔ دا په استقرايي ډول داسې تعريفېږي ($!A^N$ د کلاسیک
!!{formula}~$!A$ ``شهودي'' ژباړه وګڼئ):
\begin{align*}
!A^N & \ident \lnot\lnot !A \quad \text{د اټومي !!{formula}s $!A$ دپاره} \\
(!A \land !B)^N & \ident (!A^N \land !B^N) \\
(!A \lor !B)^N & \ident  \lnot\lnot (!A^N \lor !B^N) \\
(!A \lif !B)^N & \ident (!A^N \lif !B^N) \\
(\lforall[x][!A])^N & \ident \lforall[x][!A^N] \\
(\lexists[x][!A])^N & \ident \lnot\lnot\lexists[x][!A^N]
\end{align*}
کولموګوروف او ګليوېنکو د بياني منطق دپاره د دې ژباړې بڼې لرلې؛ د
پريديکاتي منطق دپاره ګوډل او ګېنڅن په خپلواکه توګه وړاندې کړه۔ لرو:

\begin{thm}
\begin{enumerate}
\item $!A \liff !A^N$ په کلاسیک ډول د اثبات وړ دے۔
\item که $!A$ په کلاسیک ډول د اثبات وړ وي، $!A^N$ په شهودي ډول د
  اثبات وړ دے۔
\end{enumerate}
\end{thm}

اوس لاندې خبرې اترې تصورولے شو۔ کلاسیک رياضيپوه: ``ما $!A$ ثابت
کړے!'' شهودي رياضيپوه: ``ستا ثبوت معتبر نۀ دے۔ تا په اصل کښې
$!A^N$ ثابت کړے دے۔'' کلاسیک رياضيپوه: ``ماته سمه ده!'' د کلاسیک
رياضيپوه له نظره شهودي رياضيپوه يوازې په وړو توپيرونو نښتے، ځکه دواړه
سره برابر دي۔ خو شهودي رياضيپوه ټينګار کوي چې توپير شته۔

پام وکړئ چې پاسنۍ ژباړه يوازې له سوچه منطق سره کار لري؛ دا پوښتنه نۀ
څېړي چې د کلاسیکو او شهودي رياضياتو دپاره مناسب \emph{غيرمنطقي}
بديهي اصول کوم دي يا د دواړو تر منځ څه اړيکه ده۔ خو د پاسنۍ قضيې
لاندې لږ غځونه څه ګټور معلومات راکوي:

\begin{thm}
که $\Gamma$ په کلاسیک ډول $!A$ ثابتوي، $\Gamma^N$ په شهودي ډول
$!A^N$ ثابتوي۔
\end{thm}

په نورو ټکو، که $!A$ له ځينو فرضيو څخه په کلاسیک ډول د اثبات وړ وي،
$!A^N$ د هغو له دوه‌ګونې نفي ژباړو څخه د اثبات وړ دے۔

د دې د ښودلو دپاره چې يوه جمله يا بياني !!{formula} په شهودي ډول
معتبر دے، يوازې يو ثبوت ورکول پکار دي۔ خو دا به څنګه ښيئ چې معتبر نۀ
دے؟ د دې دپاره داسې معنٰی پوهنې ته اړ يو چې سمه، او غوره دا چې بشپړه
وي۔ د کرېپکي معنٰی پوهنه دا اړتيا ښه پوره کوي۔

هماغه کار کولے شو چې د کلاسیک منطق دپاره مو وکړ: معنٰی پوهنه تعريف
کړو او سموالے او بشپړتيا ثابته کړو۔ خو لاندې توپير ته پام ارزښت لري۔
د کلاسیک منطق معنٰی پوهنه په يوه معنا ``څرګنده'' وه او د تړونکو په
معنا کښې ضمني پرته وه۔ که څه هم د کرېپکي د معنٰی پوهنې دپاره څه
وجداني انګېزه ورکولے شو، دا د هماغه ډول حتمي والي احساس نۀ راکوي۔
سربېره پر دې، د يوه کلاسیک !!{structure} مفهوم يو طبيعي رياضيکي
مفهوم دے؛ نو يا د !!a{structure} مفهوم د کلاسیک د لومړۍ درجې منطق د
مطالعې وسيله ګڼلے شو، يا کلاسیک د لومړۍ درجې منطق د رياضيکي
!!{structure}s د مطالعې وسيله۔ د دې پر خلاف، د کرېپکي !!{structure}s
يوازې منطقي جوړونې ګڼل کېدے شي؛ داسې نۀ ښکاري چې خپلواک رياضيکي ارزښت
ولري۔

د بياني ژبې دپاره د کرېپکي !!{structure}~$\mModel M = \tuple{W, R, V}$
د يوه سټ~$W$، پر~$W$ د يوه جزوي ترتيب~$R$ چې تر ټولو کوچنے
!!{element} ولري، او د بياني متغيرونو له يوې ``يکنواختې'' ګومارنې
څخه د~$W$ !!{element}s ته جوړ دے۔ وجداني تصور دا دے چې د~$W$
!!{element}s ``نړۍګانې'' يا ``د پوهې حالتونه'' ښيي؛ غړے
$v \geq u$ د~$u$ يو ``ممکن راتلونکے حالت'' ښيي؛ او~$u$ ته ګومارل شوي
بياني متغيرونه هغه قضيې دي چې په حالت~$u$ کښې رښتونې پېژندل شوې دي۔
بيا د تحميل اړيکه~$\mSat{M}{!A}[w]$ دا اړيکه د ژبې ټولو
!!{formula}s ته غځوي؛ $\mSat{M}{!A}[w]$ داسې ولولئ چې ``$!A$ په
حالت~$w$ کښې رښتونے دے''۔ اړيکه په استقرايي ډول داسې تعريفېږي:
\begin{enumerate}
\item $\mSat{M}{\Obj p_i}[w]$ هله او يوازې هله چې $\Obj p_i$ يو له
  هغو بياني متغيرونو څخه وي چې~$w$ ته ګومارل شوي۔
\item $\mSat/{M}{\lfalse}[w]$۔
\item $\mSat{M}{(!A \land !B)}[w]$ هله او يوازې هله چې
  $\mSat{M}{!A}[w]$ او $\mSat{M}{!B}[w]$ وي۔
\item $\mSat{M}{(!A \lor !B)}[w]$ هله او يوازې هله چې
  $\mSat{M}{!A}[w]$ يا $\mSat{M}{!B}[w]$ وي۔
\item $\mSat{M}{(!A \lif !B)}[w]$ هله او يوازې هله چې هر کله
  $w' \geq w$ او $\mSat{M}{!A}[w']$ وي، نو $\mSat{M}{!B}[w']$ هم وي۔
\end{enumerate}
دا ښودل يو ښه تمرين دے چې
$\lnot (p \land q) \lif (\lnot p \lor \lnot q)$ په شهودي ډول معتبر
نۀ دے؛ د کرېپکي داسې !!{structure} جوړ کړئ چې مقابله بېلګه ورکړي۔

\end{document}
